% =============================================================================
% NOTICE CODICOLOGIQUE, STEMMATIQUE ET PALÉOGRAPHIQUE SUR LE CODEX RAVENNAS 429
% Pour l'édition critique d'Aristophane, Les Thesmophories (Louis Hauseux, « Bibi »)
% Folios 162v à 177r (Classense 164v à 179r) — Vers 1 à 1231
% =============================================================================

\section*{Le \emph{Codex Ravennas 429 (olim 137, 4, A)} : codicologie, paléographie et établissement du texte des \emph{Thesmophories}}

\subsection*{1. Fiche codicologique, origine constantinopolitaine et histoire du manuscrit ($\mathbf{R}$)}

Le \textbf{\emph{Codex Ravennas 429}} (ancienne cote camaldule : \textbf{\emph{137, 4, A}}), conservé à Ravenne à l'\emph{Istituzione Biblioteca Classense}, est le plus ancien et le plus complet de tous les manuscrits médiévaux d'Aristophane, et le \textbf{seul témoin médiéval indépendant (\emph{codex unicus})} des \emph{Thesmophories}.

\begin{itemize}
  \item \textbf{Date et origine :} Copié à \textbf{Constantinople vers le milieu du X\textsuperscript{e} siècle} (ca. 950 apr. J.-C.), sous le règne de Constantin~VII Porphyrogénète, dans l'atelier ou l'entourage immédiat du célèbre moine et calligraphe \textbf{Éphrem} (\foreignlanguage{greek}{Ἐφραίμ μοναχός}), à qui l'on doit notamment le \emph{Venetus A} de l'\emph{Iliade} (\emph{Marcianus gr.~454}), l'Aristote \emph{Marcianus gr.~201} (954) et le \emph{Vaticanus gr.~124} de Polybe (947). L'écriture de la première main ($\text{R}^1$) et celle du réviseur-scholiaste ($\text{R}^2$) présentent une parenté jumelle avec celles du célèbre \emph{Laurentianus Plut.~32.9} (Sophocle, Eschyle, Apollonios de Rhodes).
  \item \textbf{Description matérielle et double foliotation :} Codex de parchemin de \textbf{191~folios} (numérotés de \texttt{1} à \texttt{191} dans le fac-similé phototypique de J.~van~Leeuwen, 1904, mais de \texttt{3} à \texttt{193} au crayon dans l'angle supérieur droit par la Biblioteca Classense, qui inclut deux feuillets de garde initiaux : $\text{folio van Leeuwen} + 2 = \text{folio Classense}$). Dimensions : env. $322 \times 210\text{ mm}$ (justification de la colonne poétique : env. $225 \times 110\text{ mm}$), composé de quaternions réguliers réglés à la pointe sèche. Le texte poétique est disposé sur une seule colonne à raison de \textbf{38 à 45 lignes par page} (généralement 42 ou 43 lignes dans les \emph{Thesmophories}), entouré sur trois côtés d'un cadre destiné aux scholies ($\Sigma^{\text{R}}$).
  \item \textbf{Contenu :} Seul manuscrit médiéval à réunir les \textbf{onze comédies conservées} d'Aristophane (\emph{Ploutos}, \emph{Nuées}, \emph{Grenouilles}, \emph{Oiseaux}, \emph{Cavaliers}, \emph{Paix}, \emph{Lysistrata}, \emph{Acharniens}, \emph{Guêpes}, \emph{Thesmophories}, \emph{Assemblée des femmes}). \textbf{Les \emph{Thesmophories} occupent les folios \texttt{162v} à \texttt{177r} (Classense \texttt{164v} à \texttt{179r}), soit exactement 30 pages (15 planches doubles, vues IIIF \texttt{166} à \texttt{180}) pour 1231 vers.}
  \item \textbf{Histoire moderne, apographe $\mathbf{M}$ et édition Juntine (1516) :} Conservé à la bibliothèque ducale d'Urbino, puis transporté à Pise au monastère camaldule de San Michele in Borgo en 1515 par Antonio Surian, il est copié en Italie pour donner naissance au \textbf{\emph{Monacensis gr.~492}} ($\mathbf{M}$, \texttt{\textbackslash Mon}, jadis \emph{Augustanus}, noté $\text{G}$ par von Velsen). Cet apographe $\text{M}$ sert à l'humaniste florentin \textbf{Eufrosino Bonini} (\emph{Euphrosynus Boninus}) pour l'\emph{editio princeps} de \emph{Lysistrata} et des \emph{Thesmophories} publiée à Florence chez Bernardo Giunta le 28 janvier 1516 (ces deux comédies ayant fait défaut à l'Aldine de Marcus Musurus en 1498). Racheté en Toscane en 1712 par l'abbé camaldule \textbf{Pietro Canneti} pour le monastère de Classe à Ravenne (cote \texttt{137, 4, A}), le manuscrit est publié par Filippo Invernizi (1794), minutieusement collationné par Adolf von Velsen (1871, 1883) et pour les scholies par Albert Martin (1882), avant d'être reproduit en phototypie par Johannes van Leeuwen (1904).
\end{itemize}

\subsection*{2. Témoins papyrologiques ($\mathbf{\Psi^1}$ à $\mathbf{\Psi^5}$), tradition indirecte et «~Dossier Mnésiloque~»}

\subsubsection*{A. Les cinq papyri antiques confrontés à $\mathbf{R}$}
Conformément au \emph{Conspectus siglorum} de l'édition (\texttt{Aristophane-Thesmo Introduction.tex}, ll.~371--380), cinq témoins papyrologiques permettent de contrôler le texte de $\text{R}$ :
\begin{enumerate}
  \item \textbf{$\mathbf{\Psi^1}$ (\texttt{\textbackslash PsiUn}, \emph{PSI}~XI~1194 + XIV p.~XV, II\textsuperscript{e}~s.) :} couvre les \textbf{vv.~133--144, 145--156, 237--245, 273--288, 594--596, 804--809, 1231} (\texttt{ff.~164r, 165r, 165v, 169v, 172r, 177r}). Il confirme au v.~141 \foreignlanguage{greek}{σύ τ'} (avec $\Sigma$) contre \foreignlanguage{greek}{τίς δ'} de $\text{R}$ ; restitue au v.~242 \foreignlanguage{greek}{[ἓ]τερον} (Medaglia) omis par $\text{R}$ ; atteste après le v.~276 la \emph{parepigraphè} \foreignlanguage{greek}{ολολυ[} ; donne au v.~285 le pluriel \foreignlanguage{greek}{τὰ πόπαν'} (conjecturé par Porson) contre \foreignlanguage{greek}{τὸ πόπανον} de $\text{R}$ ; au v.~287 la forme attique \foreignlanguage{greek}{Φερρέφαττα} contre \foreignlanguage{greek}{Φερσέφαττα} de $\text{R}$ ; au v.~807 \foreignlanguage{greek}{Στρατονίκην} contre \foreignlanguage{greek}{στρατωνίκην} de $\text{R}$ ; et conserve après le v.~1231 la stichométrie alexandrine de \textbf{1201 vers}.
  \item \textbf{$\mathbf{\Psi^2}$ (\texttt{\textbackslash PsiDeux}, \emph{P.Oxy.}~LVI~3839, éd. W.~E.~H.~Cockle, ca.~200) :} couvre les \textbf{vv.~742--766 (fr.~1) et 941--953 (fr.~2)} (\texttt{ff.~171r, 171v, 173v}). Il marque au v.~743 le changement d'interlocuteur avant \foreignlanguage{greek}{εἰπέ μοι} par un point haut et un blanc de deux lettres ; confirme au v.~745 la géminée \foreignlanguage{greek}{τυννοῦτον / τυννοῦτο} (contre \foreignlanguage{greek}{τυνοῦτον / τυνοῦτο} de $\text{R}$) et l'\textbf{absence de changement de personnage} avant \foreignlanguage{greek}{μικρὸν} (contre le sigle fautif \foreignlanguage{greek}{Γϋ.} de $\text{R}^2$) ; donne au v.~746 \foreignlanguage{greek}{γέγονε} et au v.~754 \foreignlanguage{greek}{δός τὸ σφαγεῖον} (sans le \foreignlanguage{greek}{μοι} parasite de $\text{R}$).
  \item \textbf{$\mathbf{\Psi^3}$ (\texttt{\textbackslash PsiTrois}, \emph{P.Oxy.}~LVI~3840, éd. P.~J.~Parsons, IV\textsuperscript{e}~s.) :} couvre les \textbf{vv.~1185--1193} (\texttt{f.~176v}). Il confirme au v.~1185 le nominatif barbare \foreignlanguage{greek}{γογγυλί} (conjecturé par Enger contre \foreignlanguage{greek}{γογγύλῃ} de $\text{R}$), conserve après le v.~1187 sur une ligne distincte la \emph{parepigraphè} \foreignlanguage{greek}{μ]ενο[} (\foreignlanguage{greek}{[ἀπεψωλη]μένο[ς]}), et confirme au v.~1192 \foreignlanguage{greek}{μέλι} (avec $\text{M}$, contre \foreignlanguage{greek}{μέλις} de $\text{R}$).
  \item \textbf{$\mathbf{\Psi^4}$ (\texttt{\textbackslash PsiQuatre}, \emph{P.Oxy.}~LXXIII~4935, éd. A.~Benaissa + LXXVIII~5132, éd. W.~B.~Henry, II\textsuperscript{e}~s.) :} couvre les \textbf{vv.~1043--1051 (fr.~1) et 1201--1210 (fr.~2)} (\texttt{ff.~174v, 175r, 176v}). Il confirme au v.~1047 \foreignlanguage{greek}{[ι]ω μοιρ[ας} sans le \foreignlanguage{greek}{μοι} dittographique de $\text{R}$ (confirmant Hermann) ; confirme au v.~1051 \foreignlanguage{greek}{βάρβαρον} de $\text{R}$ (défendu par Alphonse Willems et Louis Hauseux) ; et montre au v.~1208 un espace initial trop court pour \foreignlanguage{greek}{λέλυσον} ($\text{R}^{\text{ac}}$), appuyant \foreignlanguage{greek}{λέλυσο} de $\text{R}^{\text{pc}}$.
  \item \textbf{$\mathbf{\Psi^5}$ (\texttt{\textbackslash PsiCinq}, \emph{P.Oxy.}~IX~1176, fr.~39 col.~XII, Satyrus, \emph{Vie d'Euripide}, II\textsuperscript{e}~s.) :} cite les \textbf{vv.~335--337} (\foreignlanguage{greek}{ευριπιδιωι[..]υτ} face à \foreignlanguage{greek}{Εὐριπίδῃ Μήδοις} de $\text{R}$) et \textbf{vv.~373--375} (donnant la variante \foreignlanguage{greek}{Ἀρ[χίκλ]ει'} suppl. Hunt face à \foreignlanguage{greek}{Τιμόκλει'} de $\text{R}$, $\Sigma^{\text{R}}$ et la Souda $\varepsilon~2616$).
\end{enumerate}

\subsubsection*{B. «~Le dossier Mnésiloque~» : authenticité du nom \foreignlanguage{greek}{Μνησίλοχος} transmis par $\mathbf{R^2}$ (\texttt{f.~162v})}
En tête du \texttt{f.~162v}, l'\emph{Hypothesis} et la liste des personnages copiées par $\text{R}^2$ nomment le beau-père d'Euripide \textbf{\foreignlanguage{greek}{Μνησίλοχος}}. Contre le scepticisme des éditeurs modernes qui lui substituent l'anonyme \foreignlanguage{greek}{Κηδεστής}, Louis Hauseux (\texttt{Introduction.tex}, ll.~325--357) démontre l'authenticité historique et dramatique de ce nom :
\begin{enumerate}
  \item La \emph{Vie d'Euripide} atteste que l'épouse du tragique, \textbf{Choirilè} (ou \textbf{Choirinè}), était fille de \textbf{Mnésiloque}, et que le deuxième fils d'Euripide portait le nom de \textbf{Mnésiloque} --- conformément à l'usage athénien constant (cf. Démosthène, 43.74) voulant que le cadet porte le nom de son grand-père maternel (comme l'aîné, Mnésarchidès, portait celui du père d'Euripide).
  \item Au \textbf{v.~289 (\texttt{f.~165v})}, dans sa prière aux Deux Déesses pour ses enfants, là où $\text{R}$ porte \foreignlanguage{greek}{χοῖρον} et $\Sigma^{\text{R}}$ \foreignlanguage{greek}{χοιρίον}, \textbf{\texttt{Bibi} restitue le nom propre \foreignlanguage{greek}{Χοιρίλ'} (ou \foreignlanguage{greek}{Χοιρίν'})}, c'est-à-dire précisément Choirilè/Choirinè, fille de Mnésiloque et épouse d'Euripide, à qui son père souhaite malicieusement de trouver un mari riche et sot, tout le contraire de son gendre Euripide.
\end{enumerate}

\subsection*{3. Paléographie des trois mains ($\mathbf{R^1}$, $\mathbf{R^2}$ et Euphrosynus Boninus) et guide des abréviations}

L'examen autoptique des planches \texttt{ff.~162v--177r} conduit à distinguer trois mains :
\begin{enumerate}
  \item \textbf{La première main ($\mathbf{R^1}$ / \texttt{\textbackslash Rac}, milieu du X\textsuperscript{e}~s.) :} copiste du titre et de la colonne poétique des vers~1 à 1231 (ainsi que des quatre \emph{parepigraphai} intégrées à la colonne). Écriture minuscule livresque penchée à tendance «~bouletée~». $\text{R}^1$ n'écrit jamais les noms des locuteurs : il signale les changements de réplique uniquement par la \textbf{\emph{paragraphos}} (\texttt{—}) sous le début du vers et, à l'intérieur du vers, par le \textbf{\emph{dikolon}} (\texttt{:}). Le iota muet est toujours adscrit lorsqu'il est noté (\foreignlanguage{greek}{τῶι, ῆι, τουτωί} v.~250).
  \item \textbf{La seconde main ($\mathbf{R^2}$ / \texttt{\textbackslash Rdeux}, milieu du X\textsuperscript{e}~s.) :} réviseur et scholiaste contemporain (\emph{scholiorum manus, aequalis fere primae manui}, von Velsen). Utilisant une semi-onciale de glose (\emph{Kleinoziale}) riche en abréviations tachygraphiques par superposition (\texttt{\textbackslash auDessus}), $\text{R}^2$ a copié l'\emph{Hypothesis} et l'index des personnages au \texttt{f.~162v}, les scholies marginales et interlinéaires ($\Sigma^{\text{R}}$), les sigles de personnages (\foreignlanguage{greek}{μνησ/, κηδ/, ευρ/, αγαθ/, χορ/, γυν/, κλεισθ/, πρυτ/, τοξ/}), les quatre notes scéniques marginales de déguisement aux \texttt{ff.~175r--175v} (v.~1056 \foreignlanguage{greek}{ευρ}$^{\pi}_{\iota}$ \foreignlanguage{greek}{ηχω}, v.~1065 \foreignlanguage{greek}{ὁ μνησίλοχος ὡς ἀνδρομέδα}, v.~1098 \foreignlanguage{greek}{ευρ}$^{\pi}_{\iota}$ \foreignlanguage{greek}{ως περσευς}, v.~1107 \foreignlanguage{greek}{κ}$^{\delta}_{\eta}$ \foreignlanguage{greek}{ως ανδρομ}$^{\delta}_{\varepsilon}$), et a effectué de nombreuses corrections sur grattage (\emph{in rasura}, \texttt{\textbackslash inRas}) ainsi que l'ajout en marge des vers omis par $\text{R}^1$ (\textbf{v.~33} au \texttt{f.~162v} et \textbf{vv.~1102--1103} dans la marge inférieure du \texttt{f.~175v}).
  \item \textbf{La troisième main humaniste (Euphrosynus Boninus, 1515/1516) :} comme l'ont établi von Velsen (1871) et Albert Martin (1882, p.~V) --- rappelés par Louis Hauseux aux vv.~1059 n.~5 et 1197--1216 ---, l'éditeur de la Juntine \textbf{Eufrosino Bonini} a inscrit de sa propre main quelques notes marginales de locuteur sur le \emph{Ravennas}, notamment l'abréviation \textbf{\foreignlanguage{greek}{γϋ.}} (= \foreignlanguage{greek}{γυνή}, désignant Euripide déguisé en vieille maquerelle) aux \textbf{vv.~1189 et 1198 (\texttt{f.~176v})}.
\end{enumerate}

\subsection*{4. Les quatre didascalies scéniques antiques (\emph{parepigraphai}) conservées par $\mathbf{R}$}

\begin{enumerate}
  \item \textbf{Post v.~99 (\texttt{f.~163v}) :} \foreignlanguage{greek}{μινυρισμός} (\emph{«~fredonnement vocal d'Agathon~»}).
  \item \textbf{Post v.~129 (\texttt{f.~164r}) :} \foreignlanguage{greek}{ὀλολύζει ὁ γέρων} (\emph{«~le vieillard [Mnésiloque] pousse le cri rituel d'ololyguè~»}).
  \item \textbf{Post v.~276 (\texttt{f.~165v}) :} \foreignlanguage{greek}{ὀλολύζουσι τε ἱερὸν ὠθεῖται} (confirmé par $\Psi^1$ \foreignlanguage{greek}{ολολυ[}, accompagné d'une \emph{koronis} marginale et de la scholie $\Sigma^{\text{R}}$ \foreignlanguage{greek}{παρεπιγραφή. ἐκκυκλεῖται ἐπὶ τὸ ἔξω τὸ θεσμοφόριον}).
  \item \textbf{Post v.~1187 (\texttt{f.~176v}) :} \foreignlanguage{greek}{ἀνακύπτη· καὶ παρακύπτι· ἀπεψωλημένος·} (= \foreignlanguage{greek}{ἀνακύπτει καὶ παρακύπτει ἀπεψωλημένος}, confirmé sur une ligne distincte par $\Psi^3$ \foreignlanguage{greek}{μ]ενο[}, supprimé du texte parlé par Ellebodius).
\end{enumerate}

\subsection*{5. Typologie des confusions paléographiques et d'abréviations élucidées par \texttt{Bibi}}

\subsubsection*{A. Confusions propres à la minuscule bouletée et aux signes tachygraphiques}
\begin{itemize}
  \item \textbf{$\nu \leftrightarrow \mu$, $\iota \leftrightarrow \varsigma$ et abréviation tachygraphique de \foreignlanguage{greek}{προ-} (en forme de $\text{Ʒ}$ inversé semblable à $\varepsilon$, cf. Gardthausen II p.~340) :}
  \begin{itemize}
    \item Au \textbf{v.~50 (\texttt{f.~163r})}, $\text{R}$ porte \foreignlanguage{greek}{πράμος ἡμέτερος}. \textbf{\texttt{Bibi} restitue \foreignlanguage{greek}{πρανοῦν προτέρους}} (d'après Hésychius $\pi~3188$ et \emph{Nuées}~562) : \foreignlanguage{greek}{πρανοι ν} médivisé a donné \foreignlanguage{greek}{πράμος} ($\nu \to \mu$, $\iota \to \varsigma$), tandis que \foreignlanguage{greek}{προτέρους} abrégé par le signe tachygraphique de \foreignlanguage{greek}{προ-} a été rattaché au $\nu$ précédent et lu \foreignlanguage{greek}{<η>μ-ετερος}.
    \item Aux \textbf{vv.~162--163 (\texttt{f.~164r})}, $\text{R}$ porte \foreignlanguage{greek}{ἐχύμησαν ... διεκίνων}. \textbf{\texttt{Bibi} restitue \foreignlanguage{greek}{ἐκεχήνεσαν ... διέκειντ'}} : haplographie \foreignlanguage{greek}{εκεχ-} $\to$ \foreignlanguage{greek}{εχ-}, iotacisme $\eta \to \upsilon$ et confusion minuscule $\nu \to \mu$.
    \item Au \textbf{v.~710 (\texttt{f.~170v})}, $\text{R}$ porte \foreignlanguage{greek}{ἥκεις τ'} : \textbf{\texttt{Bibi} restitue \foreignlanguage{greek}{ἥκειν γ'}} d'après Photius $\eta~88$ (permutation minuscule $\nu \to \varsigma$).
  \end{itemize}
  \item \textbf{Bêta minuscule en forme de $u$ à deux pointes confondu avec kappa ($\kappa \leftrightarrow \beta$) :}
  \begin{itemize}
    \item Au \textbf{v.~895 (\texttt{f.~173r})}, $\text{R}$ porte \foreignlanguage{greek}{βαϋζα} : \textbf{\texttt{Bibi} restitue \foreignlanguage{greek}{κόνυζα}} («~l'inule visqueuse~», plante rituelle des Thesmophories).
    \item Dans la \textbf{scholie $\Sigma^{\text{R}}$ au v.~1174 (\texttt{f.~176v})}, $\text{R}$ porte \foreignlanguage{greek}{τὸ ἁβρῶς βαδίζειν} : \textbf{\texttt{Bibi} restitue \foreignlanguage{greek}{τὸ ἄκρως βαδίζειν}} d'après Zonaras p.~1171.
  \end{itemize}
  \item \textbf{Haplographie, prothèse ou métathèse de $\sigma$ :}
  \begin{itemize}
    \item Au \textbf{v.~556 (\texttt{f.~169r})}, après \foreignlanguage{greek}{ὡς}, $\text{R}$ écrit \textbf{\foreignlanguage{greek}{τλεγγίδας} sans $\sigma$ initial} (corrigé en \foreignlanguage{greek}{στλεγγίδας} par $\text{M}$).
    \item Au \textbf{v.~557 (\texttt{f.~169r})}, là où $\text{R}$, $\Sigma$ et la Souda donnent \foreignlanguage{greek}{τὸν σῖτον}, \textbf{\texttt{Bibi} restitue \foreignlanguage{greek}{τὸν κύσθον}} par métathèse du $\sigma$ et iotacisme $\upsilon/\iota$, exactement parallèle au \textbf{v.~1114 (\texttt{f.~175v})} où $\text{R}$ écrit \foreignlanguage{greek}{σκύτο} pour \foreignlanguage{greek}{κύστο} (= \foreignlanguage{greek}{κύσθον}) et au \textbf{v.~1188 (\texttt{f.~176v})} où $\text{R}^{\text{ac}}$ avait écrit \foreignlanguage{greek}{τὸ σποστίον} (corrigé en \foreignlanguage{greek}{τὸ }\texttt{\textbackslash inRas\{σ\}\{\}}\foreignlanguage{greek}{ποστίον} par $\text{R}^{\text{pc}}$).
    \item Au \textbf{v.~834 (\texttt{f.~172r})}, après \foreignlanguage{greek}{δίδοσθαι}, $\text{R}$ écrit \foreignlanguage{greek}{Τηνίοισι} par chute du $\Sigma$ initial, là où $\Sigma^{\text{R}}$ et la Souda ($\sigma~1089$) conservent \textbf{\foreignlanguage{greek}{Στηνίοισι}} (la fête des \emph{Stènia}).
  \end{itemize}
\end{itemize}

\subsubsection*{B. Confusions remontant à l'archétype en onciale (\emph{scriptio continua})}
\begin{itemize}
  \item \textbf{Abréviation onciale $\text{Π'} = \text{περ}$ confondue avec $\text{Γ'} = \text{τρίτη}$ :} au \textbf{v.~80 (\texttt{f.~163v})}, $\text{R}$ porte \foreignlanguage{greek}{ἐπεὶ τρίτη 'στὶ}, issu de \textbf{\foreignlanguage{greek}{ἐπείπερ ἐστὶ}} (Nauck, adopté par \texttt{Bibi} : \texttt{ΕΠΕΙΠ'ΕϹΤΙ} lu \texttt{ΕΠΕΙΓ'ΕϹΤΙ}).
  \item \textbf{Chute de $\Lambda$ devant $\text{Α}$ en onciale (\texttt{ΛΑ} $\to$ \texttt{Α}) :} au \textbf{v.~910 (\texttt{f.~173r})}, \textbf{$\text{R}$ porte \foreignlanguage{greek}{ἀπὸ ἀφύων}} ($\Sigma$ et Souda \foreignlanguage{greek}{ἰφύων}). \textbf{\texttt{Bibi} restitue \foreignlanguage{greek}{ἀπὸ λαιφίων}} («~au vu de tes haillons~», d'après Hésychius $\lambda~171$ \foreignlanguage{greek}{λαίφια· ῥάκη}) : dans \texttt{ΑΠΟΛΑΙΦΙΩΝ}, le $\Lambda$ triangulaire placé devant $\text{Α}$ est tombé par haplographie (\texttt{ΑΠΟΑΙΦΙΩΝ}), puis l'iotacisme a produit \foreignlanguage{greek}{ἀφύων}.
  \item \textbf{Confusion onciale $\text{ΛΙ} \leftrightarrow \text{ΑΝ}$ :} au \textbf{v.~1040 (\texttt{f.~174v})}, $\text{R}$ porte \foreignlanguage{greek}{λιτομέναν} (\texttt{ΛΙΤΟΜΕΝΑΝ}, \foreignlanguage{greek}{λιτομένα} Enger) face à la variante \foreignlanguage{greek}{ἀντομένα(ν)} (\texttt{ΑΝΤΟΜΕΝΑΝ}) de $\Sigma^{\text{R}}$.
  \item \textbf{Confusion onciale $\text{ΙϹ} \leftrightarrow \text{Π}$ :} au \textbf{v.~657 (\texttt{f.~170r})}, $\text{R}$ donne \foreignlanguage{greek}{εἰσελήλυθεν} (\texttt{ΕΙϹΕΛΗΛΥΘΕΝ}), corrigé par \textbf{Handley} en \textbf{\foreignlanguage{greek}{ἐπελήλυθε}} (\texttt{ΕΠΕΛΗΛΥΘΕ}).
  \item \textbf{Confusion onciale $\text{Τ} \leftrightarrow \text{Γ}$ et intrusion de glose :} au \textbf{v.~702 (\texttt{f.~170v})}, $\text{R}$ donne \foreignlanguage{greek}{ὡς ἅπαν γὰρ ἐστὶν τόλμης ἔργα κἀναισχυντίας} pour \textbf{\foreignlanguage{greek}{ὡς ἅπαντ' ἄρ' ἐστὶ τόλμης πλέα κἀναισχυντίας}} (Bentley + Souda : \texttt{ΑΠΑΝΤΑΡ} lu \texttt{ΑΠΑΝΓΑΡ}, et glose \foreignlanguage{greek}{ἔργα} remplaçant \foreignlanguage{greek}{πλέα} ou \foreignlanguage{greek}{διπλὰ} abrégé avec $\ddot{\delta}$).
  \item \textbf{Confusion onciale $\text{Κ} \leftrightarrow \text{Λ}$ :} au \textbf{v.~711 (\texttt{f.~170v})}, $\text{R}$ et la Souda donnent \foreignlanguage{greek}{οὐ λέξεις} ; appuyé sur la scholie interne de $\text{R}$ (\foreignlanguage{greek}{ὅθεν οὐ φεύξῃ}), \textbf{\texttt{Bibi} restitue \foreignlanguage{greek}{οὐκ ἔξει.}} (\texttt{ΟΥΚΕΞΕΙ} lu \texttt{ΟΥΛΕΞΕΙ[Σ]}).
  \item \textbf{Confusion des rondes $\text{Ο} \leftrightarrow \text{Θ}$ :} au \textbf{v.~815 (\texttt{f.~172r})}, l'attique \foreignlanguage{greek}{ποοῦντας} (\texttt{ΠΟΟΥΝΤΑϹ}) est devenu \foreignlanguage{greek}{ποθοῦντας} (\texttt{ΠΟΘΟΥΝΤΑϹ}) dans $\text{R}$ (\foreignlanguage{greek}{ποιοῦντας} Zanetti).
  \item \textbf{Médivisions de \emph{scriptio continua} :} \textbf{v.~781 (\texttt{f.~171v})} \foreignlanguage{greek}{τὸ ῥῶ} $\text{R} \to$ \textbf{\foreignlanguage{greek}{τορῶ} \texttt{Bibi}} (ap. Hésychius $\tau~1162$) ; \textbf{v.~839 (\texttt{f.~172v})} \foreignlanguage{greek}{ὦ πόλεις} $\text{R} \to$ \textbf{\foreignlanguage{greek}{ωὖπολι} \texttt{Bibi}} (adresse à Eupolis, datant la pièce de 410 av. J.-C.) ; \textbf{v.~1071 (\texttt{f.~175r})} \foreignlanguage{greek}{μέρος ἐξέλαχον;} $\text{R} \to$ \textbf{\foreignlanguage{greek}{μόρος ἔξελ' ἄχος.} \texttt{Bibi}}.
\end{itemize}

\subsection*{6. Tableau synoptique des 30 folios des \emph{Thesmophories} dans le \emph{Ravennas}}

\begin{center}
\small
\begin{tabular}{|c|c|c|c|l|}
\hline
\textbf{N\textsuperscript{o}} & \textbf{Folio (Leeuwen)} & \textbf{Classense} & \textbf{Planche IIIF} & \textbf{Vers contenus et particularités majeures} \\
\hline
01 & \texttt{f.~162v} & \texttt{164v} & 166 (g.) & \emph{Hypothesis} ($\text{R}^2$) + Titre + \textbf{vv.~1--42} (v.~33 in mg. $\text{R}^{\text{pc}}$, om. $\text{M}$) \\
02 & \texttt{f.~163r} & \texttt{165r} & 166 (dr.) & \textbf{vv.~43--79} (v.~45 \emph{in ras.} ; v.~50 \foreignlanguage{greek}{πρανοῦν προτέρους} \texttt{Bibi}) \\
03 & \texttt{f.~163v} & \texttt{165v} & 167 (g.) & \textbf{vv.~80--122} (v.~80 Nauck ; post v.~99 \emph{parep.} \foreignlanguage{greek}{μινυρισμός}) \\
04 & \texttt{f.~164r} & \texttt{166r} & 167 (dr.) & \textbf{vv.~123--164} (post v.~129 \foreignlanguage{greek}{ὀλολύζει ὁ γέρων} ; $\Psi^1$ vv.~133--156 ; vv.~162--163 \texttt{Bibi}) \\
05 & \texttt{f.~164v} & \texttt{166v} & 168 (g.) & \textbf{vv.~165--207} \\
06 & \texttt{f.~165r} & \texttt{167r} & 168 (dr.) & \textbf{vv.~208--251} ($\Psi^1$ vv.~237--245 ; v.~242 \foreignlanguage{greek}{[ἓ]τερον} $\Psi^1$ om. $\text{R}$) \\
07 & \texttt{f.~165v} & \texttt{167v} & 169 (g.) & \textbf{vv.~252--294} ($\Psi^1$ vv.~273--288 ; post v.~276 \emph{parep.} ; vv.~289, 291 \texttt{Bibi}) \\
08 & \texttt{f.~166r} & \texttt{168r} & 169 (dr.) & \textbf{vv.~295--336} (gloses vv.~297--298 ; v.~320 \emph{in ras.} ; $\Psi^5$ vv.~335--336) \\
09 & \texttt{f.~166v} & \texttt{168v} & 170 (g.) & \textbf{vv.~337--379} ($\Psi^5$ vv.~337, 373--375 ; v.~354 Willems) \\
10 & \texttt{f.~167r} & \texttt{169r} & 170 (dr.) & \textbf{vv.~380--422} (v.~383 \foreignlanguage{greek}{καλλιλεξία γυνή} $\text{R}^{\text{pc}}$) \\
11 & \texttt{f.~167v} & \texttt{169v} & 171 (g.) & \textbf{vv.~423--465} (v.~446 \foreignlanguage{greek}{ἔγκυος} \texttt{Bibi}) \\
12 & \texttt{f.~168r} & \texttt{170r} & 171 (dr.) & \textbf{vv.~466--508} (v.~477 \texttt{Bibi} ; vv.~489, 493, 506 $\text{R}^{\text{pc}}$) \\
13 & \texttt{f.~168v} & \texttt{170v} & 172 (g.) & \textbf{vv.~509--550} (v.~512 \emph{in ras.} $\text{R}^{\text{pc}}$) \\
14 & \texttt{f.~169r} & \texttt{171r} & 172 (dr.) & \textbf{vv.~551--589} (v.~556 \foreignlanguage{greek}{τλεγγίδας} $\text{R}$ ; v.~557 \foreignlanguage{greek}{τὸν κύσθον} \texttt{Bibi}) \\
15 & \texttt{f.~169v} & \texttt{171v} & 173 (g.) & \textbf{vv.~590--634} ($\Psi^1$ vv.~594--596 ; v.~630 \foreignlanguage{greek}{μέντοι} in mg. $\text{R}^{\text{pc}}$, om. $\text{M}$) \\
16 & \texttt{f.~170r} & \texttt{172r} & 173 (dr.) & \textbf{vv.~635--678} (vv.~637--643 \texttt{Bibi} ; v.~657 Handley ; v.~669 \emph{in ras.}) \\
17 & \texttt{f.~170v} & \texttt{172v} & 174 (g.) & \textbf{vv.~679--720} (v.~702 ; vv.~704, 710, 711 \texttt{Bibi}) \\
18 & \texttt{f.~171r} & \texttt{173r} & 174 (dr.) & \textbf{vv.~721--764} (v.~721 \texttt{Bibi} ; $\Psi^2$ vv.~742--764) \\
19 & \texttt{f.~171v} & \texttt{173v} & 175 (g.) & \textbf{vv.~765--799} ($\Psi^2$ vv.~765--766 ; v.~781 \foreignlanguage{greek}{τορῶ} \texttt{Bibi} ; début Parabase) \\
20 & \texttt{f.~172r} & \texttt{174r} & 175 (dr.) & \textbf{vv.~800--837} ($\Psi^1$ vv.~804--809 ; vv.~822--823 \texttt{Bibi} ; v.~834 \foreignlanguage{greek}{Στηνίοισι}) \\
21 & \texttt{f.~172v} & \texttt{174v} & 176 (g.) & \textbf{vv.~838--875} (v.~839 \foreignlanguage{greek}{ωὖπολι} \texttt{Bibi} ; v.~853 \foreignlanguage{greek}{ἑλένην} \texttt{Bibi}) \\
22 & \texttt{f.~173r} & \texttt{175r} & 176 (dr.) & \textbf{vv.~876--919} (v.~895 \foreignlanguage{greek}{κόνυζα} \texttt{Bibi} ; v.~910 \foreignlanguage{greek}{λαιφίων} \texttt{Bibi} pro \foreignlanguage{greek}{ἀφύων} $\text{R}$) \\
23 & \texttt{f.~173v} & \texttt{175v} & 177 (g.) & \textbf{vv.~920--958} ($\Psi^2$ vv.~941--953 ; v.~942 \emph{in ras.} $\text{R}^2$) \\
24 & \texttt{f.~174r} & \texttt{176r} & 177 (dr.) & \textbf{vv.~959--1000} (ode chorale circulaire) \\
25 & \texttt{f.~174v} & \texttt{176v} & 178 (g.) & \textbf{vv.~1001--1043} (v.~1006 \emph{in ras.} ; vv.~1031, 1041 \texttt{Bibi} ; $\Psi^4$ v.~1043) \\
26 & \texttt{f.~175r} & \texttt{177r} & 178 (dr.) & \textbf{vv.~1044--1081} ($\Psi^4$ vv.~1044--1051 ; notes mg. vv.~1056, 1065 ; vv.~1071--1079 \texttt{Bibi}) \\
27 & \texttt{f.~175v} & \texttt{177v} & 179 (g.) & \textbf{vv.~1082--1127} (v.~1086 \emph{in ras.} ; vv.~1090--1093 \texttt{Bibi} ; vv.~1102--1103 in mg. inf. $\text{R}^2$) \\
28 & \texttt{f.~176r} & \texttt{178r} & 179 (dr.) & \textbf{vv.~1128--1168} (post v.~1143 del. $\text{R}^{\text{pc}}$ ; v.~1151 \foreignlanguage{greek}{θεμιτεύετον} \texttt{Bibi}) \\
29 & \texttt{f.~176v} & \texttt{178v} & 180 (g.) & \textbf{vv.~1169--1210} ($\Psi^3$, $\Psi^4$ ; post v.~1187 \emph{parep.} ; 6 \emph{in ras.} ; vv.~1197--1198 \texttt{Bibi}) \\
30 & \texttt{f.~177r} & \texttt{179r} & 180 (dr.) & \textbf{vv.~1211--1231 (Fin)} ($\Psi^1$ v.~1231 ; v.~1215 \foreignlanguage{greek}{συμβήνη} \texttt{Bibi} ; + \emph{Hyp. Eccl.}) \\
\hline
\end{tabular}
\end{center}

\subsection*{7. Notices analytiques, paléographiques et critiques folio par folio (\texttt{ff.~162v--177r})}

\begin{description}
  \item[\textbf{Folio 01 --- \texttt{f.~162v} (\texttt{Classense 164v}, planche 166 g.) : \emph{Hypothesis} + vv.~1--42.}]
  En haut de page, $\text{R}^1$ a réservé 12 lignes remplies par $\text{R}^2$ en semi-onciale de glose à abréviations tachygraphiques (\emph{Hypothesis} et liste des personnages donnant le nom \foreignlanguage{greek}{Μνησίλοχος}). Titre de première main : \foreignlanguage{greek}{※ ἀριστοφάνους θεσμοφοριάζουσαι —}.
  \emph{Interventions de $\text{R}^{\text{pc}}$ :} v.~28 copié après v.~30 par $\text{R}^{\text{ac}}$, rétabli par $\text{R}^{\text{pc}}$ ; v.~32 \foreignlanguage{greek}{(Μνη.)} $\text{R}^{\text{ac}} \to$ \foreignlanguage{greek}{(Εὐρ.) continuat} $\text{R}^{\text{pc}}, \Sigma$ (\foreignlanguage{greek}{ἑόρακας} Bentley pro \foreignlanguage{greek}{ἑόρακα} $\text{R}$) ; \textbf{v.~33 omis dans la colonne par $\text{R}^{\text{ac}}$ et par $\text{M}$, ajouté en marge par $\text{R}^{\text{pc}}$} ; v.~34 \foreignlanguage{greek}{καὶ μεγ'} $\text{R}^{\text{ac}}$ (\foreignlanguage{greek}{ὶ} eras. $\text{R}^{\text{pc}}$, \foreignlanguage{greek}{καὶ μεγ'} $\text{M}$, \foreignlanguage{greek}{κἀμέ γ'} edd.).
  \emph{Leçons et choix de l'édition :} v.~3 \foreignlanguage{greek}{τι} $\text{R}$ : \foreignlanguage{greek}{τε} $\Sigma$, Souda ; v.~10 \foreignlanguage{greek}{ὁρᾶν.} $\text{R}$ (maintenu par \texttt{Bibi}) ; v.~18 \foreignlanguage{greek}{ἀκοὴν δὲ χοάνης} $\text{R}$, Willems ; v.~21 \foreignlanguage{greek}{τε} $\text{R}$ : \foreignlanguage{greek}{γε} $\Sigma$ ; v.~23 del. Willems ; v.~24 \foreignlanguage{greek}{προσμάθοι· μὴ} $\text{R}$ : \foreignlanguage{greek}{προσμάθοιμι} Ellebodius ; \textbf{vv.~30--31 \foreignlanguage{greek}{Ἁγάθων} et \foreignlanguage{greek}{ἀγαθῶν} (\foreignlanguage{greek}{(Μνη.) continuat} $\text{R}$) \texttt{Bibi}} ; v.~39 \foreignlanguage{greek}{λέως} $\text{R}$, Souda AF : \foreignlanguage{greek}{λαός} $\Sigma$, Souda GIVM.

  \item[\textbf{Folio 02 --- \texttt{f.~163r} (\texttt{Classense 165r}, planche 166 dr.) : vv.~43--79.}]
  \emph{Interventions de $\text{R}^{\text{pc}}$ :} v.~45 \foreignlanguage{greek}{λέγει}\texttt{\textbackslash inRas\{ς\}\{\}} $\text{R}^{\text{pc}}$ ; v.~60 \foreignlanguage{greek}{κατὰ} om. $\text{R}^{\text{ac}}$, add. $\text{R}^{\text{pc}}$.
  \emph{Leçons et conjectures de \texttt{Bibi} :} v.~44 \foreignlanguage{greek}{δὲ} $\text{R}$ ; \textbf{v.~50 \foreignlanguage{greek}{πράμος ἡμέτερος} $\text{R} \to$ \foreignlanguage{greek}{πρανοῦν προτέρους} \texttt{Bibi}} (ap. Hésychius $\pi~3188$ et \emph{Nub.}~562 : médivision de \foreignlanguage{greek}{πρανοι ν} avec $\nu \to \mu, \iota \to \varsigma$, et signe tachygraphique de \foreignlanguage{greek}{προ-} en forme de $\text{Ʒ}$ inversé lu $\varepsilon$) ; v.~53 \foreignlanguage{greek}{ἀσπίδας} $\text{R}$ : \foreignlanguage{greek}{ἁψῖδας} Souda $\delta~1547, \lambda~181$ ; v.~56 \foreignlanguage{greek}{γογγύλίζει} $\text{R}$, Souda : \foreignlanguage{greek}{γογγύλλει} Porson ; v.~61 \foreignlanguage{greek}{γογγυλίσας} $\text{R}$ : \foreignlanguage{greek}{συγγογγύλας} Enger ; v.~69 \foreignlanguage{greek}{θύρασι} $\text{R}$ ; \textbf{vv.~71--74 attribués intégralement à \foreignlanguage{greek}{(Μνη.)} par \texttt{Bibi} ap. $\Sigma$} (contre $\text{R}$ qui donnait 72--74 à \foreignlanguage{greek}{Εὐρ.}) ; v.~74 \foreignlanguage{greek}{χρῆν ... ἐμόν} $\text{R}$.

  \item[\textbf{Folio 03 --- \texttt{f.~163v} (\texttt{Classense 165v}, planche 167 g.) : vv.~80--122.}]
  Contient après le v.~99 la \textbf{1\textsuperscript{re} \emph{parepigraphè}} : \foreignlanguage{greek}{μινυρισμός}. Dans la monodie d'Agathon (vv.~101--122), $\text{R}^2$ introduit à tort en marge l'alternance \foreignlanguage{greek}{αγαθ/ — χορ/}.
  \emph{Leçons et conjectures :} \textbf{v.~80 \foreignlanguage{greek}{ἐπεὶ τρίτη 'στὶ} $\text{R}$ : \foreignlanguage{greek}{ἐπείπερ ἐστὶ} Nauck} (abréviation onciale $\text{Π'} = \text{περ}$ lue $\text{Γ'} = \text{τρίτη}$) ; v.~86 \foreignlanguage{greek}{δί' και'} $\text{R}$ : \foreignlanguage{greek}{δικαίαν} Wilamowitz ; v.~96 \foreignlanguage{greek}{ποῖο ἐστὶν οὗτος — οὑκκυκλούμενος} $\text{R}$ : \foreignlanguage{greek}{ποῖος ἐστὶν; — οὗτος οὑκκυκλούμενος} Tyrwhitt ; v.~99 \foreignlanguage{greek}{ἂν} $\text{R}$ : \foreignlanguage{greek}{νῦν} Rogers ; \textbf{v.~100 \foreignlanguage{greek}{Μύρμηκος} \texttt{Bibi}} ap. Hésychius $\mu~1904$ ; v.~101 \foreignlanguage{greek}{χθονίαις} $\text{R}$ : \foreignlanguage{greek}{χθονίαιν} Meineke ; v.~105 \foreignlanguage{greek}{νιν ... εὐπείστως} Dindorf, Reiske ; v.~107 \foreignlanguage{greek}{ὅπλιζε} $\text{R}$ : \foreignlanguage{greek}{ὄλβιζε} Bentley ; v.~115 \foreignlanguage{greek}{ἀεισαντ'} $\text{R}$ : \foreignlanguage{greek}{ἀείσατ'} Zanetti ; v.~122 \foreignlanguage{greek}{διανεύματα} $\text{R}$ : \foreignlanguage{greek}{διὰ νεύματα} Fritzsche.

  \item[\textbf{Folio 04 --- \texttt{f.~164r} (\texttt{Classense 166r}, planche 167 dr.) : vv.~123--164 ($\mathbf{\Psi^1}$ vv.~133--156).}]
  Contient après le v.~129 la \textbf{2\textsuperscript{e} \emph{parepigraphè}} : \foreignlanguage{greek}{ὀλολύζει ὁ γέρων}.
  \emph{Leçons et conjectures :} v.~126 \foreignlanguage{greek}{φῶς} $\text{R}$ : \foreignlanguage{greek}{φάος} Burges ; v.~128 \foreignlanguage{greek}{Φοῖβον τιμᾷ} $\text{R}$ (\foreignlanguage{greek}{τιμᾷ} glose intruse) : \foreignlanguage{greek}{Φοῖβον} Dindorf ; v.~130 \foreignlanguage{greek}{πότνια} $\text{R}$ : \foreignlanguage{greek}{πότνιαι} Souda ; v.~134 \foreignlanguage{greek}{νεανίσχ' εἴ τις} $\text{R}$ : \foreignlanguage{greek}{νεανίσκ' εἴ τις} $\text{M}$ ; \textbf{v.~141 \foreignlanguage{greek}{σύ τ'} $\Psi^1, \Sigma$ : \foreignlanguage{greek}{τίς δ'} $\text{R}$} ; v.~154 \foreignlanguage{greek}{ἄν} $\Psi^1, \text{R}$ ; \textbf{vv.~162--163 \foreignlanguage{greek}{ἐχύμησαν ... διεκίνων} $\text{R} \to$ \foreignlanguage{greek}{ἐκεχήνεσαν ... διέκειντ'} \texttt{Bibi}} (haplographie \foreignlanguage{greek}{εκεχ} $\to$ \foreignlanguage{greek}{εχ}, iotacisme $\eta/\upsilon$ et confusion minuscule $\nu/\mu$).

  \item[\textbf{Folio 05 --- \texttt{f.~164v} (\texttt{Classense 166v}, planche 168 g.) : vv.~165--207.}]
  vv.~168--169 \foreignlanguage{greek}{Φιλοκλῆς, Ξενοκλῆς} $\text{R}$ : \foreignlanguage{greek}{Φιλοκλέης, Ξενοκλέης} Bentley ; v.~173 \foreignlanguage{greek}{ἦν} $\text{R}$, Souda : \foreignlanguage{greek}{ἦ} Dindorf ; v.~176 \foreignlanguage{greek}{(Ἀγάθ.) λέγε} $\text{R}$ ; v.~178 \foreignlanguage{greek}{οἷόν} $\text{R}$ : \foreignlanguage{greek}{οἷός} Souda, Stobée ; v.~179 \foreignlanguage{greek}{κοινῇ} $\text{R}$, Souda : \foreignlanguage{greek}{καινῇ} Biset ap. $\Sigma$ ; v.~196 \foreignlanguage{greek}{γὰρ} $\text{R}$ : \foreignlanguage{greek}{γὰρ ἂν} Souda $\upsilon~709$ ; v.~204 \foreignlanguage{greek}{νυκτερείσια} $\text{R}$.

  \item[\textbf{Folio 06 --- \texttt{f.~165r} (\texttt{Classense 167r}, planche 168 dr.) : vv.~208--251 ($\mathbf{\Psi^1}$ vv.~237--245).}]
  v.~216 \foreignlanguage{greek}{ἀφεύειν} $\text{R}$ | \foreignlanguage{greek}{είσω} $\text{R}$ : \foreignlanguage{greek}{εἴ σοι} Bonini ; \textbf{v.~217 \foreignlanguage{greek}{διδόναι γ' αὐτὸν} $\text{R} \to$ \foreignlanguage{greek}{διδόναι γ' ἑαυτὸν} \texttt{Bibi} auctore Willems} ; v.~220 \foreignlanguage{greek}{ξυροδίκης} $\text{R}$ : \foreignlanguage{greek}{ξυροδόκης} $\Sigma$, Pollux ; v.~223 \foreignlanguage{greek}{ἀττατα· ατταται} $\text{R}$ : \foreignlanguage{greek}{ἀτταταῖ ἰατταταῖ} Biset ; v.~234 \foreignlanguage{greek}{θεάσασθαι σεαυτόν} $\text{R}$ : \foreignlanguage{greek}{θεάσασθ' αὑτόν} Fritzsche ; \textbf{v.~242 \foreignlanguage{greek}{[ἓ]τερον} $\Psi^1$ (Medaglia) : om. $\text{R}$} ; v.~248 \foreignlanguage{greek}{οἰμώζετ' ... εἰς} $\text{R}$ : \foreignlanguage{greek}{οἰμώξετ' ... εἴ τις} Zanetti, Brunck ; v.~250 \foreignlanguage{greek}{τουτωί} $\text{R}$ (iota adscrit).

  \item[\textbf{Folio 07 --- \texttt{f.~165v} (\texttt{Classense 167v}, planche 169 g.) : vv.~252--294 ($\mathbf{\Psi^1}$ vv.~273--288).}]
  Contient après le v.~276 la \textbf{3\textsuperscript{e} \emph{parepigraphè}} : \foreignlanguage{greek}{ὀλολύζουσι τε ἱερὸν ὠθεῖται} $\text{R}$ (\foreignlanguage{greek}{ολολυ[} $\Psi^1$, avec \emph{koronis} et $\Sigma^{\text{R}}$ \foreignlanguage{greek}{παρεπιγραφή. ἐκκυκλεῖται ἐπὶ τὸ ἔξω τὸ θεσμοφόριον}).
  \emph{Leçons et conjectures de \texttt{Bibi} :} vv.~252--260 attributions de $\text{R}^2$ corrigées par Bentley, Dobree, Ellebodius ; \textbf{vv.~263--264 \foreignlanguage{greek}{(Ἀγάθ.)} maintenu au v.~263 (\foreignlanguage{greek}{γ' ο}$^{\tilde{\upsilon}}_{\nu}$ $\text{R}^{\text{pc}}$) et suppléé au v.~264 par \texttt{Bibi}} ; v.~273 \foreignlanguage{greek}{Ὑποκράτους} $\text{R}$ (calembour maintenu par \texttt{Bibi}) ; \textbf{v.~285 \foreignlanguage{greek}{τὰ πόπαν'} $\Psi^1$, Porson : \foreignlanguage{greek}{τὸ πόπανον} $\text{R}$} ; \textbf{v.~287 \foreignlanguage{greek}{Φερρέφαττα} $\Psi^1$, Thiersch : \foreignlanguage{greek}{Φερσέφαττα} $\text{R}$} ; \textbf{v.~289 \foreignlanguage{greek}{χοῖρον} $\text{R}$ (\foreignlanguage{greek}{χοιρίον} $\Sigma$) $\to$ \foreignlanguage{greek}{Χοιρίλ'} \texttt{Bibi}} (fille de Mnésiloque et épouse d'Euripide) ; \textbf{v.~291 \foreignlanguage{greek}{πρὸς θάληκον} $\text{R} \to$ \foreignlanguage{greek}{πρὸς θυλάκιον} (\foreignlanguage{greek}{πρὸς θάλικα} ap. Hésychius $\theta~50$) \texttt{Bibi}}.

  \item[\textbf{Folio 08 --- \texttt{f.~166r} (\texttt{Classense 168r}, planche 169 dr.) : vv.~295--336 ($\mathbf{\Psi^5}$ vv.~335--336).}]
  vv.~297--298 gloses intruses dans $\text{R}$ (\foreignlanguage{greek}{τῇ Δήμητρῃ καὶ τῇ Κόρῃ} del. Blaydes ; \foreignlanguage{greek}{τῇ Γῇ} del. Dobree ex $\Sigma$) ; v.~306 \texttt{\textbackslash expo\{καὶ\}}\foreignlanguage{greek}{τὴν ἀγορεύουσαν} $\text{R}^2$ ; \textbf{v.~320 \foreignlanguage{greek}{θηροφ}\texttt{\textbackslash inRas\{ω\}\{ό\}}\foreignlanguage{greek}{νε παῖ--} $\text{R}^{\text{pc}}$} (\foreignlanguage{greek}{θηροφόνη} Hermann) ; v.~329 \foreignlanguage{greek}{Ἀθηναίων} $\text{R}$ : \foreignlanguage{greek}{Ἀθηνῶν} Reisig ; v.~331 \foreignlanguage{greek}{(Κηρ.)} $\text{R}$ : \foreignlanguage{greek}{(Κορ.)} van Leeuwen.

  \item[\textbf{Folio 09 --- \texttt{f.~166v} (\texttt{Classense 168v}, planche 170 g.) : vv.~337--379 ($\mathbf{\Psi^5}$ vv.~337, 373--375).}]
  v.~337 \foreignlanguage{greek}{ευριπιδιωι[..]υτ} $\Psi^5$ : \foreignlanguage{greek}{Εὐριπίδῃ Μήδοις} $\text{R}$ (\foreignlanguage{greek}{Μήδοις τ'} Biset) ; \textbf{v.~354 \foreignlanguage{greek}{γενέσθαι} $\text{R} \to$ \foreignlanguage{greek}{ἀποτελεῖσθαι} Willems} ; \textbf{v.~374 \foreignlanguage{greek}{Ἀρ[χίκλ]ει'} $\Psi^5$ (Hunt) : \foreignlanguage{greek}{Τιμόκλει'} $\text{R}, \Sigma$, Souda $\varepsilon~2616$}.

  \item[\textbf{Folio 10 --- \texttt{f.~167r} (\texttt{Classense 169r}, planche 170 dr.) : vv.~380--422.}]
  \textbf{v.~383 \foreignlanguage{greek}{(}\texttt{\textbackslash auDessus\{γυ\}\{ν\}}\foreignlanguage{greek}{)} $\text{R}^{\text{ac}} \to$ \foreignlanguage{greek}{(καλλιλεξία γυνή)} $\text{R}^{\text{pc}}, \text{M}$} (\foreignlanguage{greek}{(Μίκα)} Rogers) ; vv.~392--393 \foreignlanguage{greek}{μυχοτρόπους, ἀνδρεαστρίας, οἰνοπότιδας} $\text{R} \to$ \foreignlanguage{greek}{μοιχοτρόπους, ἀνδρεραστρίας, οἰνοπίπας} Souda ; v.~400 \foreignlanguage{greek}{τις πλέκῃ} $\text{R} \to$ \foreignlanguage{greek}{τις ξυμπλέκῃ} Willems ; v.~417 \foreignlanguage{greek}{μυχοῖς} $\text{R} \to$ \foreignlanguage{greek}{μοιχοῖς} $\text{M}$ ; v.~419 \foreignlanguage{greek}{ταμιεύεσθαι} $\text{R} \to$ \foreignlanguage{greek}{ταμιεῦσαι καὶ} Reiske.

  \item[\textbf{Folio 11 --- \texttt{f.~167v} (\texttt{Classense 169v}, planche 171 g.) : vv.~423--465.}]
  v.~433 \foreignlanguage{greek}{οὔπω τε} $\text{R}^{\text{ac}} \to$ \foreignlanguage{greek}{οὐπώποτε} $\text{R}^{\text{pc}}$ (\foreignlanguage{greek}{οὔπω} Dindorf) ; v.~437 \foreignlanguage{greek}{εἰδέας} $\text{R}$ (\foreignlanguage{greek}{ἰδέας} Souda) ; v.~443 \foreignlanguage{greek}{γυνὴ ἑτέρα} $\text{R}^{\text{pc}}$ ; \textbf{v.~446 \foreignlanguage{greek}{ἐν Κύπρῳ} $\text{R} \to$ \foreignlanguage{greek}{ἔγκυος} \texttt{Bibi}} ; v.~456 \foreignlanguage{greek}{τοῖς} Plutarque, Aulu-Gelle : om. $\text{R}$ ; v.~462 \foreignlanguage{greek}{ἄκερα} $\text{R} \to$ \foreignlanguage{greek}{ἄκαιρα} Zanetti.

  \item[\textbf{Folio 12 --- \texttt{f.~168r} (\texttt{Classense 170r}, planche 171 dr.) : vv.~466--508.}]
  \textbf{v.~477 \foreignlanguage{greek}{πολλ'} $\text{R} \to$ \foreignlanguage{greek}{πολλά γ'· ἀλλ'} \texttt{Bibi}} ; \textbf{v.~489 \foreignlanguage{greek}{κυνδ'} $\text{R}^{\text{ac}} \to$ \foreignlanguage{greek}{κύβδ'} $\text{R}^{\text{pc}}, \text{M}$} ; \textbf{v.~493 \foreignlanguage{greek}{ληκώμεθα} $\text{R}^{\text{ac}}$ (?), Souda $\to$ \foreignlanguage{greek}{κινώμεθα} $\text{R}^{\text{pc}}$} ; v.~500 del. van Herwerden ; v.~501 \foreignlanguage{greek}{μυχὸν} $\text{R} \to$ \foreignlanguage{greek}{μοιχὸν} Grynaeus ; \textbf{v.~506 \foreignlanguage{greek}{βοῴη} suppl. $\text{R}^{\text{pc}}$}.

  \item[\textbf{Folio 13 --- \texttt{f.~168v} (\texttt{Classense 170v}, planche 172 g.) : vv.~509--550.}]
  \textbf{v.~512 \texttt{\textbackslash auDessus\{εἰ\}\{η\}\textbackslash inRas\{σ\}\{\}}\foreignlanguage{greek}{έφερε} $\text{R}^{\text{pc}}$} (\foreignlanguage{greek}{ἥ 'φερεν} Bentley) ; v.~514 \foreignlanguage{greek}{αὖτ' ἔκμαγμα} (ex \foreignlanguage{greek}{έγμαγμα}) $\text{R}^{\text{pc}}$ (\foreignlanguage{greek}{αὐτέκμαγμα} Reiske) ; v.~533 \foreignlanguage{greek}{(γυνὴ τις)} $\text{R}^{\text{pc}}$ | \foreignlanguage{greek}{Ἄγραυλον} $\text{R} \to$ \foreignlanguage{greek}{Ἄγλαυρον} Daubuz ; v.~549 \foreignlanguage{greek}{ἂν} ex \foreignlanguage{greek}{ἂρ} corr. $\text{R}^{\text{pc}}$ ; v.~550 \foreignlanguage{greek}{δ' ἁπαξαπάσας} $\text{M}^{\text{ac}}$.

  \item[\textbf{Folio 14 --- \texttt{f.~169r} (\texttt{Classense 171r}, planche 172 dr.) : vv.~551--589.}]
  \textbf{v.~556 \foreignlanguage{greek}{τλεγγίδας} $\text{R}$ (haplographie de $\sigma$ après \foreignlanguage{greek}{ὡς}) $\to$ \foreignlanguage{greek}{στλεγγίδας} $\text{M}$} ; \textbf{v.~557 \foreignlanguage{greek}{τὸν σῖτον} $\text{R}, \Sigma$, Souda $\to$ \foreignlanguage{greek}{τὸν κύσθον} \texttt{Bibi}} (cf. v.~1114) ; v.~569 \foreignlanguage{greek}{πρόσθες} $\text{R} \to$ \foreignlanguage{greek}{πρόσθιγε} Willems ; v.~581 \foreignlanguage{greek}{ἀφράκτοις} $\text{R} \to$ \foreignlanguage{greek}{ἀφάρκτοις} Dindorf.

  \item[\textbf{Folio 15 --- \texttt{f.~169v} (\texttt{Classense 171v}, planche 173 g.) : vv.~590--634 ($\mathbf{\Psi^1}$ vv.~594--596).}]
  v.~596 \foreignlanguage{greek}{πεπύσμην} $\text{R}, \Psi^1$ (\foreignlanguage{greek}{'πεπύσμην} Faber) ; \textbf{v.~630 \foreignlanguage{greek}{μέντοι} $\text{R}^{\text{pc}}$ in mg., Souda $\pi~2563$ : om. $\text{R}^{\text{ac}}, \text{M}$} ; \textbf{v.~634 \foreignlanguage{greek}{δεῦρο} $\text{R}^{\text{ac}}, \text{M} \to$ \foreignlanguage{greek}{δευρ ὦ} $\text{R}^2$} | \foreignlanguage{greek}{Κλεισόσθενες} $\text{R} \to$ \foreignlanguage{greek}{Κλείσθενες} Bonini.

  \item[\textbf{Folio 16 --- \texttt{f.~170r} (\texttt{Classense 172r}, planche 173 dr.) : vv.~635--678.}]
  \textbf{vv.~637--643 redistribués par \texttt{Bibi} : \foreignlanguage{greek}{(Μνη.)} 637--638 (avec \foreignlanguage{greek}{χαλᾷ} v.~638), \foreignlanguage{greek}{(Μίκα)} 639--640, \foreignlanguage{greek}{(Κλεισθ.)} 642--643} ; \textbf{v.~657 \foreignlanguage{greek}{εἰσελήλυθεν} $\text{R} \to$ \foreignlanguage{greek}{ἐπελήλυθε} Handley} (confusion onciale $\text{ΙϹ/Π}$) ; v.~659 \foreignlanguage{greek}{(Ἡμιχ. γυν.)} $\text{R}^{\text{pc}}$ ; \textbf{v.~669 \foreignlanguage{greek}{ἅπασιν }\texttt{\textbackslash inRas\{ἅπασιν\}\{\}}\foreignlanguage{greek}{ ἔσται} $\text{R}^{\text{pc}}$} (\foreignlanguage{greek}{ἀνδράσιν ἔσται} Beer ap. Hermann).

  \item[\textbf{Folio 17 --- \texttt{f.~170v} (\texttt{Classense 172v}, planche 174 g.) : vv.~679--720.}]
  v.~686 \foreignlanguage{greek}{ἀποτίνεται· | παραχρῆμά τε τίνεται} $\text{R} \to$ \foreignlanguage{greek}{παραχρῆμ' ἀποτίνεται} Hermann ; v.~689 \foreignlanguage{greek}{ἀά} supra lin. $\text{R}$ ; \textbf{v.~693 \foreignlanguage{greek}{ἀφε}\texttt{\textbackslash auDessus\{ῖ\}\{ῆ\}}\foreignlanguage{greek}{τ'} $\text{R}^2$} ; \textbf{v.~702 \foreignlanguage{greek}{ἅπαν γὰρ ἐστὶν ... ἔργα} $\text{R} \to$ \foreignlanguage{greek}{ἅπαντ' ἄρ' ἐστὶ ... πλέα} Bentley + Souda} ; \textbf{v.~704 \foreignlanguage{greek}{ἐξάρξω} $\text{R} \to$ \foreignlanguage{greek}{ἐξαρήξω} \texttt{Bibi}} ap. Hésychius $\varepsilon~3599$ ; \textbf{v.~710 \foreignlanguage{greek}{ἥκεις τ'} $\text{R} \to$ \foreignlanguage{greek}{ἥκειν γ'} \texttt{Bibi}} ap. Photius $\eta~88$ ; \textbf{v.~711 \foreignlanguage{greek}{οὐ λέξεις} $\text{R}$, Souda $\to$ \foreignlanguage{greek}{οὐκ ἔξει.} \texttt{Bibi}} ap. $\Sigma^{\text{R}}$.

  \item[\textbf{Folio 18 --- \texttt{f.~171r} (\texttt{Classense 173r}, planche 174 dr.) : vv.~721--764 ($\mathbf{\Psi^2}$ fr.~1, vv.~742--764).}]
  \textbf{v.~721 \foreignlanguage{greek}{ξὺν ἀθέοις ἔργοις γὰρ} \texttt{Bibi} ap. Taillardat} ; \textbf{v.~743 changement de locuteur avant \foreignlanguage{greek}{εἰπέ μοι} dans $\Psi^2$} ; \textbf{v.~745 \foreignlanguage{greek}{τυννοῦτον / τυννοῦτο} et absence de \foreignlanguage{greek}{Γϋ.} devant \foreignlanguage{greek}{μικρὸν} dans $\Psi^2$} ; v.~746 \foreignlanguage{greek}{γέγονε} $\Psi^2$ ; \textbf{v.~754 \foreignlanguage{greek}{δός τὸ} $\Psi^2$, Lobeck (contre \foreignlanguage{greek}{δός μοι τὸ} $\text{R}$) et \foreignlanguage{greek}{σφαγίον} $\text{R}^{\text{ac}} \to$ \foreignlanguage{greek}{σφαγεῖον} $\text{R}^{\text{pc}}, \text{M}$} ; v.~759 \foreignlanguage{greek}{(ἄλλη γυνή)} $\text{R}^{\text{pc}}$ (\foreignlanguage{greek}{(Κρίτ.)} Sommerstein).

  \item[\textbf{Folio 19 --- \texttt{f.~171v} (\texttt{Classense 173v}, planche 175 g.) : vv.~765--799 ($\mathbf{\Psi^2}$ vv.~765--766).}]
  v.~772 \foreignlanguage{greek}{πόθεν; <πόθεν;>} Scaliger ; v.~773 \foreignlanguage{greek}{εἶτα διὰ τἀγάλματ'} $\text{R} \to$ \foreignlanguage{greek}{εἰ ταδὶ τἀγάλματ'} Ellebodius (\foreignlanguage{greek}{τἀγάλματα} = les \foreignlanguage{greek}{πόπανα} sacrés, Willems) ; \textbf{v.~781 \foreignlanguage{greek}{τὸ ῥῶ} $\text{R} \to$ \foreignlanguage{greek}{τορῶ} \texttt{Bibi}} ap. Hésychius $\tau~1162$ (médivision de \texttt{ΤΟΡΩ}).

  \item[\textbf{Folio 20 --- \texttt{f.~172r} (\texttt{Classense 174r}, planche 175 dr.) : vv.~800--837 ($\mathbf{\Psi^1}$ vv.~804--809).}]
  \textbf{v.~807 \foreignlanguage{greek}{Στρατονίκην} $\Psi^1$ : \foreignlanguage{greek}{στρατωνίκην} $\text{R}$} ; v.~815 \foreignlanguage{greek}{ποθοῦντας} $\text{R} \to$ \foreignlanguage{greek}{ποιοῦντας} Zanetti (de l'attique \foreignlanguage{greek}{ποοῦντας}, confusion onciale $\text{Ο/Θ}$) ; \textbf{vv.~822--823 \foreignlanguage{greek}{τἀντίον ὁ κανὼν οἱ καλαθίσκοι / τὸ σκιάδειον·} $\text{R} \to$ \foreignlanguage{greek}{ταὐτὸ κανοῦν καὶ τὸ σκιάδειον·} \texttt{Bibi}} (expulsion de la glose \foreignlanguage{greek}{οἱ καλαθίσκοι}) ; \textbf{v.~834 \foreignlanguage{greek}{Τηνίοισι} $\text{R} \to$ \foreignlanguage{greek}{Στηνίοισι} $\Sigma$, Souda $\sigma~1089$}.

  \item[\textbf{Folio 21 --- \texttt{f.~172v} (\texttt{Classense 174v}, planche 176 g.) : vv.~838--875.}]
  v.~838 \foreignlanguage{greek}{ὑστέρ}\texttt{\textbackslash auDessus\{ο\}\{α\}}\foreignlanguage{greek}{ν} $\text{R}^{\text{pc}}$ ; \textbf{v.~839 \foreignlanguage{greek}{ὦ πόλεις} $\text{R} \to$ \foreignlanguage{greek}{ωὖπολι} \texttt{Bibi}} (adresse à Eupolis, datant la pièce de 410 av. J.-C.) ; v.~844 \foreignlanguage{greek}{τ}\texttt{\textbackslash auDessus\{ὰ\}\{ο\}}\foreignlanguage{greek}{δὶ} $\text{R}^{\text{pc}}$ ; \textbf{v.~853 \foreignlanguage{greek}{Ἑλένη} $\text{R} \to$ \foreignlanguage{greek}{ἑλένην} \texttt{Bibi}} ap. Hésychius $\varepsilon~1995$ ; v.~860 \foreignlanguage{greek}{σύ} (ex \foreignlanguage{greek}{σί}) $\text{R}^2$ (\foreignlanguage{greek}{σοί} Zanetti).

  \item[\textbf{Folio 22 --- \texttt{f.~173r} (\texttt{Classense 175r}, planche 176 dr.) : vv.~876--919.}]
  v.~892 \foreignlanguage{greek}{αὐτὸν} (ex \foreignlanguage{greek}{αὖτὸν}) $\text{R}^{\text{pc}}$ (\foreignlanguage{greek}{αὖ τὸν} Casaubon) ; \textbf{v.~895 \foreignlanguage{greek}{βαϋζα} $\text{R} \to$ \foreignlanguage{greek}{κόνυζα} \texttt{Bibi}} (confusion minuscule $\kappa \to \beta$) ; \textbf{v.~910 \foreignlanguage{greek}{ἀφύων} $\text{R}$ (\foreignlanguage{greek}{ἰφύων} $\Sigma$, Souda) $\to$ \foreignlanguage{greek}{λαιφίων} \texttt{Bibi}} ap. Hésychius $\lambda~171$ (chute de $\Lambda$ devant $\text{Α}$ dans \texttt{ΑΠΟΛΑΙΦΙΩΝ}).

  \item[\textbf{Folio 23 --- \texttt{f.~173v} (\texttt{Classense 175v}, planche 177 g.) : vv.~920--958 ($\mathbf{\Psi^2}$ fr.~2, vv.~941--953).}]
  v.~928 \foreignlanguage{greek}{(Γυ.)} $\text{R} \to$ \foreignlanguage{greek}{(Μνη.)} Scaliger ; \textbf{v.~942 \texttt{\textbackslash inRas\{γε\}\{\}}\foreignlanguage{greek}{ τοῖς} $\text{R}^2$} ; v.~945 \foreignlanguage{greek}{ἰαππαπαιάξ} $\text{R}$, Souda ; v.~947 \foreignlanguage{greek}{πέσωμεν} $\text{R} \to$ \foreignlanguage{greek}{παίσωμεν} Zanetti, cf. $\Sigma$.

  \item[\textbf{Folio 24 --- \texttt{f.~174r} (\texttt{Classense 176r}, planche 177 dr.) : vv.~959--1000.}]
  Ode chorale circulaire : v.~969 \foreignlanguage{greek}{ἐλύραν} $\text{R} \to$ \foreignlanguage{greek}{εὐλύραν} Zanetti ; v.~982 \foreignlanguage{greek}{χαίρειν} $\text{R} \to$ \foreignlanguage{greek}{χάριν} Ellebodius ; v.~995 \foreignlanguage{greek}{συῒ κτυπεῖται} $\text{R} \to$ \foreignlanguage{greek}{σοὶ κτυπεῖται} Zanetti ; v.~996 \foreignlanguage{greek}{Κιθαρώνιος} $\text{R} \to$ \foreignlanguage{greek}{Κιθαιρώνιος} Zanetti.

  \item[\textbf{Folio 25 --- \texttt{f.~174v} (\texttt{Classense 176v}, planche 178 g.) : vv.~1001--1043 ($\mathbf{\Psi^4}$ v.~1043).}]
  \textbf{v.~1006 \foreignlanguage{greek}{κακοδαίμω}\texttt{\textbackslash expo\{ν\}}\foreignlanguage{greek}{ γέρ}\texttt{\textbackslash inRas\{ω\}\{ο\}}\foreignlanguage{greek}{ν} $\text{R}^{\text{pc}}, \text{M}^{\text{pc}}$} ; vv.~1015--1021 \foreignlanguage{greek}{(Εὐρ.)} $\text{R} \to$ \foreignlanguage{greek}{(Μνη.)} Tyrwhitt ; v.~1022 \foreignlanguage{greek}{(Μνη.)} praefixit $\text{R}^{\text{pc}}, \text{M}$ ; \textbf{v.~1031 \foreignlanguage{greek}{κημὸν ἕστηκ' ἔχουσ'} $\text{R} \to$ \foreignlanguage{greek}{ἕστηκ' ἔχουσα κημόν} \texttt{Bibi}} (cf. van Leeuwen 1904, pp.~XVIII--XIX) ; v.~1040 \foreignlanguage{greek}{λιτομέναν} $\text{R}$ (\foreignlanguage{greek}{ἀντομένα[ν]} $\Sigma$) $\to$ \foreignlanguage{greek}{λιτομένα} Enger ; \textbf{v.~1041 \foreignlanguage{greek}{φεύγουσαν} $\text{R} \to$ \foreignlanguage{greek}{ψεφαυγοῦς} \texttt{Bibi}} (cf. $\Sigma$ ad v.~1054 et Hésychius $\psi~135$).

  \item[\textbf{Folio 26 --- \texttt{f.~175r} (\texttt{Classense 177r}, planche 178 dr.) : vv.~1044--1081 ($\mathbf{\Psi^4}$ fr.~1, vv.~1044--1051).}]
  \textbf{v.~1047 \foreignlanguage{greek}{[ι]ω μοιρ[ας} $\Psi^4$, Hermann : \foreignlanguage{greek}{ἰώ μοι μοίρας} $\text{R}$} | \foreignlanguage{greek}{ἀνέτικτε} $\text{R} \to$ \foreignlanguage{greek}{ἄτεγκτε} Ellebodius ; \textbf{v.~1051 \foreignlanguage{greek}{βάρβαρον} $\Psi^4, \text{R}$, Willems} ; v.~1054 \foreignlanguage{greek}{δαιμόνων} $\text{R}, \Sigma \to$ \foreignlanguage{greek}{διὰ μένων} Willems ; \textbf{v.~1056 \foreignlanguage{greek}{ευρ}\texttt{\textbackslash auDessus\{ι\}\{π\}}\foreignlanguage{greek}{ ηχω} in mg. $\text{R}^{\text{pc}}, \text{M}$} ; \textbf{v.~1065 \foreignlanguage{greek}{ὁ μνησίλοχος ὡς ἀνδρομέδα} in mg. $\text{R}^{\text{pc}}, \text{M}$} ; v.~1069 \foreignlanguage{greek}{δι' Ὀλύμπου} $\text{R}$ (\foreignlanguage{greek}{†δι'† ὀλυμπίας} \texttt{Bibi} ap. Hésychius $\text{ο}~649$) ; \textbf{v.~1071 \foreignlanguage{greek}{μέρος ἐξέλαχον;} $\text{R} \to$ \foreignlanguage{greek}{μόρος ἔξελ' ἄχος.} \texttt{Bibi}} ; \textbf{v.~1074 \foreignlanguage{greek}{στωμυλωμένη} $\text{R} \to$ \foreignlanguage{greek}{στωμυλιομενής} \texttt{Bibi}} ; \textbf{v.~1078 \foreignlanguage{greek}{(Εὐρ.) 'παΰσαι} \texttt{Bibi}} ; \textbf{v.~1079 \foreignlanguage{greek}{βάλλε σύ κόρακας} \texttt{Bibi}}.

  \item[\textbf{Folio 27 --- \texttt{f.~175v} (\texttt{Classense 177v}, planche 179 g.) : vv.~1082--1127 (+ vv.~1102--1103 in mg. inf.).}]
  v.~1083 \foreignlanguage{greek}{λαλεῖς} bis $\text{R}^{\text{ac}} \to$ \foreignlanguage{greek}{λαλῖς} bis $\text{R}^2, \Sigma$ ; \textbf{v.~1086 \foreignlanguage{greek}{π}\texttt{\textbackslash inRas\{ρ\}\{\}}\foreignlanguage{greek}{ωτε} bis $\text{R}^{\text{pc}}$} (\foreignlanguage{greek}{πωτε} bis $\text{M}$) ; \textbf{vv.~1090 et 1092 attribués à \foreignlanguage{greek}{(Χορ.)} par \texttt{Bibi} ap. $\Sigma$} ; \textbf{v.~1091 \foreignlanguage{greek}{πλησίον αὕτη} $\text{R} \to$ \foreignlanguage{greek}{πλήξασ' αὕτη} \texttt{Bibi}} ; \textbf{v.~1093 \foreignlanguage{greek}{<(Εὐρ.) ποῖ ποῖ πέτει>} suppl. \texttt{Bibi}} ; \textbf{v.~1098 \foreignlanguage{greek}{ευρ}\texttt{\textbackslash auDessus\{ι\}\{π\}}\foreignlanguage{greek}{ ως περσευς} in mg. $\text{R}^{\text{pc}}, \text{M}$} ; \textbf{vv.~1102--1103 omis dans la colonne par $\text{R}^1$ et $\text{M}$, ajoutés dans la marge inférieure par $\text{R}^2$} ; \textbf{v.~1107 \foreignlanguage{greek}{κ}\texttt{\textbackslash auDessus\{η\}\{δ\}}\foreignlanguage{greek}{ ως ανδρομ}\texttt{\textbackslash auDessus\{ε\}\{δ\}} in mg. $\text{R}^{\text{pc}}, \text{M}$} ; v.~1114 \foreignlanguage{greek}{σκύτο} $\text{R} \to$ \foreignlanguage{greek}{κύστο} Scaliger ; v.~1126 \foreignlanguage{greek}{τὸ κεπαλῆς} (ex \foreignlanguage{greek}{καὶ παλῆς}) $\text{R}^{\text{pc}}, \text{M}$ (\foreignlanguage{greek}{τὸ κεπαλή σ'} Brunck) ; v.~1127 \foreignlanguage{greek}{ἀποκέκοιψ'} $\text{R}^{\text{ac}} \to$ \foreignlanguage{greek}{ἀποκεκόψοι} $\text{R}^{\text{pc}}$ (\foreignlanguage{greek}{ἀποκεκόψι} Brunck) | \foreignlanguage{greek}{πέπυσται} $\text{R} \to$ \foreignlanguage{greek}{πέπαισται} $\Sigma$.

  \item[\textbf{Folio 28 --- \texttt{f.~176r} (\texttt{Classense 178r}, planche 179 dr.) : vv.~1128--1168.}]
  v.~1135 \foreignlanguage{greek}{τῆ μαστιγᾶν} (ex \foreignlanguage{greek}{μαστηγαν}) $\text{R}^{\text{pc}}$ ; \textbf{post v.~1143 \foreignlanguage{greek}{στυγνὰς ὡσσε} $\text{R}^{\text{ac}}$ : del. $\text{R}^{\text{pc}}$ (\foreignlanguage{greek}{στυγνὰς ὡς ἐ} $\text{M}$)} ; \textbf{v.~1151 \foreignlanguage{greek}{οὗ δὴ ἀνδράσιν οὐ θεμιτὸν} $\text{R} \to$ \foreignlanguage{greek}{οὗ δὴ ἀνδράσιν οὐ θεμιτεύετον} \texttt{Bibi}} ; v.~1153 \foreignlanguage{greek}{λαμπάσι.} \texttt{Bibi}.

  \item[\textbf{Folio 29 --- \texttt{f.~176v} (\texttt{Classense 178v}, planche 180 g.) : vv.~1169--1210 ($\mathbf{\Psi^3}$ vv.~1185--1193 ; $\mathbf{\Psi^4}$ fr.~2, vv.~1201--1210).}]
  Contient après le v.~1187 la \textbf{4\textsuperscript{e} \emph{parepigraphè}} : \foreignlanguage{greek}{ἀνακύπτη· καὶ παρακύπτι· ἀπεψωλημένος·} $\text{R}$ (\foreignlanguage{greek}{μ]ενο[} $\Psi^3$, del. Ellebodius).
  \emph{Six grattages (\emph{in rasura}) et corrections :} \textbf{v.~1172 \foreignlanguage{greek}{ἐμὸν }\texttt{\textbackslash inRas\{γ'\}\{\}} $\text{R}^2$} ; v.~1173 \foreignlanguage{greek}{ἔφραζον} (ex \foreignlanguage{greek}{ἔφρον}) $\text{R}^{\text{pc}}$ ; v.~1185 \foreignlanguage{greek}{ὡ στεριπὸ} (ex \foreignlanguage{greek}{σπεριπο}) $\text{R}^{\text{pc}}$ | \textbf{\foreignlanguage{greek}{γογγυλί} $\Psi^3$, Enger : \foreignlanguage{greek}{γογγύλῃ} $\text{R}$} ; \textbf{v.~1187 \foreignlanguage{greek}{κλαυσ}\texttt{\textbackslash inRas\{' ε\}\{\}}\foreignlanguage{greek}{ϊ γ'} $\text{R}^2$} ; \textbf{v.~1188 \foreignlanguage{greek}{τὸ }\texttt{\textbackslash inRas\{σ\}\{\}}\foreignlanguage{greek}{ποστίον} $\text{R}^{\text{pc}}$} ; \textbf{v.~1192 \foreignlanguage{greek}{μέλι} $\Psi^3, \text{M}$ : \foreignlanguage{greek}{μέλις} $\text{R}$} ; \textbf{v.~1195 \texttt{\textbackslash inRas\{σὺ\}\{\}}\foreignlanguage{greek}{ κάρισοσ οὐ} $\text{R}^{\text{pc}}$} (\foreignlanguage{greek}{κάρισο σὺ} Brunck) ; \textbf{v.~1198 \foreignlanguage{greek}{κ}\texttt{\textbackslash inRas\{ω\}\{ο\}}\foreignlanguage{greek}{μίζεις αὐτοῖς} $\text{R}^{\text{pc}}$} ; \textbf{v.~1208 \foreignlanguage{greek}{λέλυσο}\texttt{\textbackslash inRas\{ν\}\{\}} $\text{R}^{\text{pc}}$} (confirmé par $\Psi^4$) ; v.~1210 \foreignlanguage{greek}{τυγάτριον} (ex \foreignlanguage{greek}{τυγάθριον}) $\text{R}^{\text{pc}}$.
  \emph{Main humaniste de Boninus :} \foreignlanguage{greek}{γϋ.} ajouté en marge aux vv.~1189 et 1198.
  \emph{Choix et conjectures de \texttt{Bibi} :} \textbf{$\Sigma^{\text{R}}$ ad v.~1174 \foreignlanguage{greek}{ἁβρῶς} $\text{R} \to$ \foreignlanguage{greek}{ἄκρως} \texttt{Bibi} ap. Zonaras p.~1171} ; \textbf{v.~1197 \foreignlanguage{greek}{συμβήνην} $\text{R}$ défendu par \texttt{Bibi} ap. Hésychius $\sigma~2137$} (contre \foreignlanguage{greek}{συβήνην} $\Sigma$, Souda, Photius) ; \textbf{v.~1198 \foreignlanguage{greek}{(Τοξότ.) κομίζι σ' αὐτίκ'} \texttt{Bibi}}.

  \item[\textbf{Folio 30 --- \texttt{f.~177r} (\texttt{Classense 179r}, planche 180 dr.) : vv.~1211--1231 (Fin, $\mathbf{\Psi^1}$ v.~1231) + \emph{Hypothesis} des \emph{Ecclesiazusae}.}]
  \textbf{v.~1215 \foreignlanguage{greek}{συβήνη} $\text{R} \to$ \foreignlanguage{greek}{συμβήνη} \texttt{Bibi}} (cf. v.~1197, Hésychius $\sigma~2137$) ; v.~1227 \foreignlanguage{greek}{πέπυσται} $\text{R} \to$ \foreignlanguage{greek}{πέπαισται} $\Sigma$ ; v.~1231 \foreignlanguage{greek}{ἀνταδοῖτον} $\text{R} \to$ \foreignlanguage{greek}{ἀνταποδοῖτον} Grynaeus (suivi dans $\Psi^1$ de la note stichométrique $1201$ vers). Dans la moitié inférieure du folio, sous la \emph{koronis}, $\text{R}^2$ copie l'\emph{Hypothesis} et l'index des personnages des \emph{Ecclesiazusae}.
\end{description}
